% ---------------------------------------------------------------------------
% DRAFT SEGMENT — HRTF measurement and the virtual HRTF set
% Standalone; not \input anywhere yet. Stitch into Methods later.
% Two \subsection{} blocks. Placed LATE in Methods, where the earmold paper put
% "Binaural recordings" and "Directional transfer functions".
% Code: hrtf/record/recordings.py::record_dome        (acquisition, head tilts)
%       hrtf/record/record_hrir.py                    (orchestration)
%       hrtf/record/processing.py                     (deconvolve, equalize,
%                                                      lowfreq, expand_azimuths)
%       hrtf/record/fit_head_radius.py                (acoustic radius fit)
%       hrtf/record/calibration/calibrate_headphones.py
% Parameters below are read off the code and the recorded params.txt, not the
% module-level interactive defaults (which differ -- head_radius in particular).
% \label{sec:hrir} resolves the placeholder reference in
% methods_donor_selection.tex.
% ---------------------------------------------------------------------------

\subsection{HRTF measurement}\label{sec:hrir}

Head-related impulse responses were measured at the blocked entrance of each
ear canal with MEMS microphones (AMM-3738-T, PUI Devices, Dayton, USA) mounted
in anthropometric earplugs of the PIRATE design
% \citep{denk_pirate_2019}
, fitted to each participant from a size series
% \citep{bau_adapting_2024}
, and supplied with phantom power (PP2B, Millenium). Logarithmic sweeps of 200\,ms duration (120\,Hz to 22\,kHz, 85\,dB SPL) were
presented ten times from each loudspeaker of the midline column and averaged
after alignment. Seven loudspeakers at $12.5^\circ$ spacing sample elevation
too coarsely to resolve the spectral cue, so the arc was recorded three times,
the participant being cued to tilt the head by one third of the loudspeaker
spacing between passes. Interleaving the physical positions in this way yields
19 elevations from $-37.5^\circ$ to $37.5^\circ$ in steps of $4.17^\circ$. A reference recording was then made with the microphones mounted
on a stand at the position the head had occupied and with no listener present,
using identical settings; it captures the response of loudspeakers,
microphones and room in situ. Emitted sweeps were pre-filtered per loudspeaker
with inverse filters obtained from a flat probe microphone, which equalizes
each loudspeaker in level and magnitude response rather than merely matching
them to one another.

Impulse responses were recovered by regularized inversion of the sweep and
each was divided by the reference response of the same loudspeaker. The
quotient is the head-related transfer function (HRTF) in the usual sense: the
pressure at the blocked ear canal entrance relative to the free-field pressure
at the position of the head with the listener absent. These are HRTFs and not
directional transfer functions: no direction-independent component has been
removed. Our earlier work
% \citep{friedrich_adaptation_2025}
reported directional transfer functions, obtained by dividing each transfer
function by the average over all measured transfer functions. That step is
omitted here because only the midline is measured, so an average over the
measured set would be an average over the very directions that are
subsequently compared, and would remove part of the elevation dependence it is
meant to preserve. Each quotient was onset-aligned, windowed with a 2.5\,ms
Hann window and truncated to 512 samples.

\subsection{Virtual HRTF set}

Because only the midline arc was measured, the remaining directions were
synthesized with a rigid-sphere model of the head
% \citep{duda_range_1998}
. The radius of the sphere was fitted for each participant to interaural time
differences measured from a horizontal row of loudspeakers, taken from the
slope of the interaural phase between 200 and 1500\,Hz. The measured arc was then replicated
across azimuth in 25 steps of $4.17^\circ$ from $-50^\circ$ to $50^\circ$,
giving 475 directions, each carrying the spherical-head magnitude shading and
interaural phase difference relative to the frontal direction at the same
elevation. Below 800\,Hz, where the measurement is limited by the sweep
bandwidth and the room, magnitude was replaced by spherical-head values.

One consequence of this construction should be stated plainly. Spectral detail
in the delivered set is measured on the midline and is identical at every
azimuth; azimuth is carried by broadband interaural level and time differences
alone, and the rendering is acoustically faithful to the measurement only at
$0^\circ$ azimuth. Three considerations make this acceptable here. The
elevation-dependent spectral pattern, which is what the manipulation alters
and what the analysis measures, is measured rather than modeled. Every
condition in the design is rendered through the identical expansion, so no
comparison rests on the synthesized component. And because a magnitude-only
edit cannot move a modeled interaural time difference, the native and modified
sets share their interaural cues exactly, which would not be guaranteed had
both been measured. The design accordingly contrasts hemifields rather than
azimuthal acuity, and no claim is made about localization in azimuth.

Two residual corrections were applied. Division by the reference leaves a
small interaural difference at frontal directions, where none should exist;
this was removed by equating broadband energy between the ears over 200\,Hz to
16\,kHz, one scalar per direction, which on the manikin left a frontal
interaural level difference of $0.000 \pm 0.003$\,dB. Headphone equalization
was measured for each participant from three re-seatings of the headphones,
averaged, and inverted as a regularized minimum-phase filter of 256 taps.

% --- open items ------------------------------------------------------------
% - CITATIONS:
%     friedrich_adaptation_2025 -- VERIFIED: Friedrich, P., and Schoenwiesner,
%       M. (2025). J. Acoust. Soc. Am. 157(3), 1543-1553,
%       doi:10.1121/10.0036056.
%     denk_pirate_2019 -- VERIFIED (Zenodo record 2574395). NOTE THE AUTHOR
%       ORDER: Denk, F., Brinkmann, F., Stirnemann, S., and Kollmeier, B.
%       (2019). "The PIRATE: an anthropometric earPlug with exchangeable
%       microphones for Individual Reliable Acquisition of Transfer functions
%       at the Ear canal entrance," Fortschritte der Akustik -- DAGA, Rostock,
%       doi:10.5281/zenodo.2574395. Denk is first author, not Brinkmann.
%     bau_adapting_2024 -- VERIFIED: Bau, D., Moscher, O., and Poerschmann, C.
%       (2024). "Adapting the PIRATE Design for Cost-Effective In-Ear
%       Microphones," Fortschritte der Akustik -- DAGA, Hannover, pp.
%       1604-1606. This is the paper the present hardware follows; it tests
%       AMM-3738-T (PUI) alongside CMM-2718AT (CUI) and SPH1642-HT5H (Knowles)
%       and finds all three near-equivalent, so it supports the choice rather
%       than prescribing it. It is also the source for the XS-XL size series.
%     Duda & Martens (1998) JASA 104(5) -- "Range dependence of the response of
%       a spherical head model". CHECK this is the right one for the model as
%       implemented; if processing.py follows a different formulation, cite
%       that instead. [NOT VERIFIED]
% - MICROPHONE now stated: AMM-3738-T (PUI Devices) + PP2B phantom supply
%   (Millenium is Thomann's house brand -- CHECK whether to name Thomann as the
%   manufacturer instead, and confirm the town for the PUI address line).
%   CONFIRM ALSO that the plugs really are the XS-XL size series rather than
%   individually printed; "fitted to each participant from a size series" is
%   written on the strength of the setup following Bau et al. (2024).
%   The old PUM-3046L-R sentence is gone -- do not let it drift back in.
% - PER-SUBJECT RADII ARE NOT REPORTED. Give the group mean and range once the
%   cohort is closed. Also state what happens on a failed fit (the code falls
%   back to 0.0875 m and stores nothing) -- if any reported participant hit the
%   fallback, that must be said.
% - The head-radius paragraph is now bare by decision (Paul, 2026-08-26): no
%   manikin validation, no Woodworth comparison. If a reviewer asks why the
%   fitted radii (~0.072 m) fall below the Woodworth 0.0875 m, the answer is
%   that the low-frequency phase delay of a real head exceeds its onset delay
%   (Kuhn 1977) and 0.0875 m is the ONSET-delay standard, inappropriate for a
%   phase-based ITD. The fit uses the same estimator the model then imposes.
%   Keep this note; it is the rebuttal, held out of the text on purpose.
% - "limited by the sweep bandwidth and the room" is my reading of why 800 Hz;
%   confirm the actual reason before submission.
% - The 19 elevations come from 7 + 6 + 6 across the three tilt passes (the
%   outermost speaker drops out of the two tilted passes). Not stated in the
%   text -- add if a reviewer queries the count.
% - record_mesm (the newer 7-speaker path) was NOT used for any reported
%   subject. If later participants are recorded on it, this subsection needs a
%   variant, and the two cohorts cannot be pooled without checking.
% - Dome loudspeaker equalization was OFF for subject and reference alike
%   (equalize_dome: False in every params.txt), which is what makes the
%   reference division a clean in-situ calibration. Not stated; probably does
%   not need to be, but it is the answer if asked why no dome EQ is mentioned.
